\documentclass[12pt]{article}
\usepackage[utf8]{inputenc}
\usepackage[T1]{fontenc}
\usepackage{hyperref}
\title{To Language Resources for K-12 Curriculum Development}
\author{Compiled by Claude Sonnet 5 (claude-sonnet-5) for Andr\'as Kornai}
\date{}
\begin{document}
\maketitle

\section{Introduction}
To (ISO 639-3 toz) is described in the search results as a language with no native speakers, used as a second language. The sources say it is a secret male initiation language of the Gbaya people, and place it in northern Cameroon and the Central African Republic. Its status is given as dormant (Verbix, Endangered Languages) or extinct (ASJP). Classification is uncertain: Glottolog places it in Atlantic-Congo, while Wikipedia lists it as unclassified in Mbum. It is unwritten.

\section{Primary Educational Resources}
None found. There is no ethnic community and no school use.

\section{Linguistic and Reference Materials}
\begin{itemize}
\item Wikipedia article on the To language. \url{https://en.wikipedia.org/wiki/To_language}
\item Glottolog entry for To (tooo1238). \url{https://glottolog.org/resource/languoid/id/tooo1238.bigmap.html}
\item ASJP word-list entry for To. \url{https://asjp.clld.org/languages/TO}
\end{itemize}

\section{Religious and Community Resources}
None found. The language has a ritual role in male initiation, which is not suitable for school materials.

\section{Digital and Audio Resources}
\begin{itemize}
\item Endangered Languages Project entry for To (North region, Mayo-Rey division, Touboro subdivision). \url{https://www.endangeredlanguages.com/lang/3517}
\end{itemize}
No audio resources were found.

\section{Curriculum Development Context}
A K-12 curriculum in To is not feasible from the resources found. There is no community of speakers and no teaching materials, and the language has a restricted ritual role. Education in the area would use French, English or the local lingua francas.
\end{document}